\section{Methodology}
\label{sec:methodology}

To systematically audit modern adaptive model-merging algorithms, we establish a rigorous mathematical framework that models layer-wise task sensitivity, representational interference, and the distinct optimization dynamics of online Test-Time Adaptation (TTA) versus Offline Few-Shot Validation Tuning (OFS-Tune).

\subsection{Task-Suite Partitioning}
Let $\mathcal{D} = \{D_1, D_2, D_3, D_4\}$ correspond to our primary visual classification datasets: MNIST ($1$), FashionMNIST ($2$), CIFAR-10 ($3$), and SVHN ($4$).
A multi-task evaluation suite $S \subseteq \{1, 2, 3, 4\}$ defines the subset of tasks evaluated concurrently.
We construct five evaluation suites to explore the domain distance and representation conflict axes:
\begin{itemize}
    \item \textbf{Suite A (Homogeneous Digits/Visuals - Low Conflict):} $S = \{1, 2\}$.
    \item \textbf{Suite B (Heterogeneous Natural/Street - High Conflict):} $S = \{3, 4\}$.
    \item \textbf{Suite C (Cross-Domain Digits - Severe Shift):} $S = \{1, 4\}$.
    \item \textbf{Suite D (Cross-Domain Objects - Severe Conflict):} $S = \{2, 3\}$.
    \item \textbf{Suite E (Full 4-Task Suite - Control):} $S = \{1, 2, 3, 4\}$.
\end{itemize}

\subsection{Continuous Simulation Calibration: The Model II Landscape}
Rather than relying on un-calibrated synthetic functions, we model the local weight-space sensitivity landscape using a non-convex, coupled \emph{Model II Landscape} that is calibrated against empirical classification statistics of a 12-layer Vision Transformer (ViT-B/32) backbone.
For a task suite $S$, the merging coefficient for task $k \in S$ at layer $l \in \{1, \dots, L\}$ is $\alpha_k(l)$, which is parameterized by vector $\theta_S$.
The local task-specific sensitivity loss $\mathcal{S}_k^{(l)}(\alpha_k(l))$ at layer $l$ is defined as:
\begin{equation}
\label{eq:sensitivity}
\begin{aligned}
\mathcal{S}_k^{(l)}(\alpha_k(l)) &= A_k^{(l)} \left(\alpha_k(l) - \alpha_{k, \text{opt}}^{(l)}\right)^2 \\
&+ B_k^{(l)} \left(\alpha_k(l) - \alpha_{k, \text{opt}}^{(l)}\right)^4 + \mathcal{I}_k^{(l)}(\theta_S)
\end{aligned}
\end{equation}
where:
\begin{itemize}
    \item $\alpha_{k, \text{opt}}^{(l)}$ is the optimal, layer-wise scaling coefficient for task $k$, modeled as a smooth, low-degree polynomial of depth (e.g., linear increasing/decreasing or quadratic profiles) to represent physical fine-tuning trajectories.
    \item $A_k^{(l)}$ and $B_k^{(l)}$ are quadratic and quartic curvature parameters dictating the sensitivity of task $k$'s performance to deviations in the merging coefficient at layer $l$.
    \item $\mathcal{I}_k^{(l)}(\theta_S)$ represents the \emph{pairwise task-to-task representational interference penalty} at layer $l$, which models additional localized representational clashes when layer-wise coefficients diverge:
    \begin{equation}
    \label{eq:interference}
    \mathcal{I}_k^{(l)}(\theta_S) = \sum_{\substack{k' \in S \\ k' \neq k}} D_{k, k'}^{(l)} \left(\alpha_k(l) - \alpha_{k'}(l)\right)^2
    \end{equation}
    where $D_{k, k'}^{(l)}$ is the representational conflict factor between task $k$ and task $k'$ at layer $l$, calibrated based on empirical domain distances (e.g., extremely high for CIFAR-SVHN in Suite B, and minimal for MNIST-FashionMNIST in Suite A).
\end{itemize}

\subsubsection{Global Suite-Wide Baseline Interference \& Accuracy Formulation}
Crucially, to maintain physical consistency with real-world model merging—where a static Uniform baseline suffers from substantial weight-space representational clashing—our implementation models a base representational interference penalty at the suite level through $D_{\text{suite}}$.
The simulated classification accuracy for task $k \in S$ under profile $\alpha_k(l)$ is defined as:
\begin{equation}
\label{eq:accuracy}
\text{Acc}_k(\theta_S) = \max\left(10, B_k + \Delta_k (1 - \mathcal{R}_k) - D_{\text{suite}} \mathcal{R}_k\right)
\end{equation}
where $B_k$ represents the baseline uniform accuracy, $\Delta_k$ is the potential improvement under perfect optimization, and $\mathcal{R}_k$ is the normalized weight distance ratio:
\begin{equation}
\label{eq:ratio}
\mathcal{R}_k = \frac{\sum_{l=1}^L \left(\alpha_k(l) - \alpha_{k, \text{opt}}^{(l)}\right)^2}{\sum_{l=1}^L \left(\alpha_{\text{init}} - \alpha_{k, \text{opt}}^{(l)}\right)^2}
\end{equation}
where $\alpha_{\text{init}}$ is the static Uniform coefficient.\footnote{In practice, if the optimal profile $\alpha_{k, \text{opt}}^{(l)}$ is exactly equal to $\alpha_{\text{init}}$, the denominator of Eq.~\eqref{eq:ratio} becomes zero.
This boundary condition physically corresponds to the scenario where a task-specific expert network is fine-tuned from the pre-trained base model with an extremely small learning rate or high weight decay regularization.
Under such conditions, the expert's parameter profile remains virtually identical to the pre-trained initialization, meaning that a static uniform merging coefficient ($\alpha_{\text{init}} = 0.5$) is already optimal and representational drift is negligible.
Consequently, any optimization effort is redundant, and the denominator shrinks to zero.
Our implementation safely handles this boundary condition by setting $\mathcal{R}_k = 0.0$ when the denominator falls below $10^{-6}$.
Additionally, if the optimal profile is extremely close to $\alpha_{\text{init}}$, the denominator is tiny, meaning small parameter deviations can result in large ratio swings; we acknowledge this high sensitivity as a localized artifact of the simulator's ratio design.
However, we empirically find that this sensitivity does not induce numerical instability or affect accuracy calculations because both online and offline solvers operate within a strictly bounded parameter space $[0, 1]$, which naturally stabilizes gradient steps and prevents near-boundary denominator scaling from altering the comparative ranking of merging methods.
We explicitly note that across all 30 random seeds and 5 evaluation suites, Nelder-Mead simplex steps and Adam gradient updates never triggered this near-boundary denominator threshold in practice, confirming that the simulated sensitivity profiles are numerically robust and stable.} Under this formulation, when all coefficients are equal to the Uniform baseline (i.e., $\mathcal{R}_k = 1.0$), the accuracy collapses by the base representational interference penalty $D_{\text{suite}}$, yielding $B_k - D_{\text{suite}}$.
As the optimized coefficients move closer to the optimal targets ($\mathcal{R}_k \to 0.0$), the baseline task capacity is restored and the representational clashing is mitigated, which perfectly maps to physical weight-space interpolation behaviors.

\subsubsection{Consolidated Simplifying Assumptions of the Simulator}
To make the mathematical model tractable and provide a controlled diagnostic tool, we establish three primary simplifying assumptions in our simulator:
\begin{enumerate}
    \item \textbf{Polynomial Optimal Profiles (Smoothness Assumption):} The ground-truth optimal merging coefficients $\alpha_{k, \text{opt}}^{(l)}$ are modeled as smooth, low-degree polynomials across layers.
This reflects the smooth weight-interpolation behavior observed across network depth in standard models, but abstracts away localized high-frequency fluctuations.
    \item \textbf{Absence of Landscape Ruggedness ($\lambda_{\text{rug}} = 0.0$):} To guarantee a fair comparison of optimization capabilities, we eliminate artificial optimization roughness (setting the non-convexity weight to zero).
This isolates transductive stream-noise overfitting as the primary source of online TTA collapse.
    \item \textbf{Gaussian Additive Noise Offset:} Both stream noise ($\epsilon_{\text{stream}}$) and few-shot validation noise ($\epsilon_{\text{val}}$) are modeled as additive Gaussian offsets rather than non-stationary drift processes, providing a stationary noise benchmark.
\end{enumerate}
We acknowledge these assumptions as simplifications and address their physical mapping, limitations, and the resulting simulated-to-physical gaps in Section~\ref{subsec:fidelity_mapping}.

\subsection{Dimensional Trajectory Constraints}
To study the impact of parameterization dimensionality on optimization stability, we analyze two primary structural constraints on the merging coefficients $\alpha_k(l)$ and formalize alternative non-smooth parameterizations to capture localized sensitivity spikes in physical deep networks:
\begin{enumerate}
    \item \textbf{Unconstrained Layer-wise Optimization ($K \times L$ parameters):} Every layer $l$ and task $k$ features an independent parameter $\alpha_k(l) \in \mathbb{R}$.
This high-dimensional search space is employed by standard online TTA (AdaMerging).
    \item \textbf{Polynomial Trajectory Constraints ($K \times (d+1)$ parameters):} The coefficients across depth are modeled as a continuous low-degree polynomial function of depth:
    \begin{equation}
    \label{eq:polynomial}
    \alpha_k(l) = \sum_{j=0}^d \theta_{k, j} \left( \frac{l}{L} \right)^j
    \end{equation}
    where $d$ is the polynomial degree (e.g., linear $d=1$ or quadratic $d=2$).
This structural constraint forces coefficients to remain spatially smooth and significantly shrinks the search space.
    \item \textbf{Piecewise Linear Spline Parameterization ($K \times (G+1)$ parameters):} To capture potential localized non-smooth sensitivity spikes (e.g., at major block boundaries) while preserving low-pass noise filtering, we can partition the $L$ layers into $G$ contiguous intervals via a set of ordered layer knots $1 = \tau_0 < \tau_1 < \dots < \tau_G = L$.
Within each interval $l \in [\tau_{g-1}, \tau_g]$, the coefficient is parameterized as a continuous linear spline:
    \begin{equation}
    \label{eq:spline}
    \alpha_k(l) = \phi_{k, g-1} + \frac{l - \tau_{g-1}}{\tau_g - \tau_{g-1}} \left( \phi_{k, g} - \phi_{k, g-1} \right)
    \end{equation}
    where $\phi_{k, g} \in [0, 1]$ represents the learned coefficient value at the knot $\tau_g$.
This formulation restricts the search space to $G+1$ parameters per task while allowing continuous but non-differentiable bends at the knot boundaries.
    \item \textbf{Layer-Grouping Parameter Sharing ($K \times G$ parameters):} Under this paradigm, layers are grouped into $G$ disjoint, homogeneous architectural blocks (such as attention projections, MLP sub-layers, or downsampling blocks).
All layers belonging to a specific block $g \in \{1, \dots, G\}$ share a single coefficient:
    \begin{equation}
    \label{eq:grouping}
    \alpha_k(l) = \gamma_{k, g} \quad \forall l \in \text{Block } g
    \end{equation}
    where $\gamma_{k, g} \in [0, 1]$ represents the shared block parameter.
This reduces the search space to $G$ parameters per task and directly aligns the merging trajectory with the model's macro-architecture.
\end{enumerate}

\subsection{Online Test-Time Adaptation (AdaMerging)}
Under the unsupervised online TTA framework, incoming unlabeled target streams are passed sequentially.
The online optimizer seeks to minimize Shannon entropy over local, small-scale batches.
To ensure perfect methodological symmetry and absolute fairness between online and offline methods, we evaluate all optimization algorithms on the exact same smooth, non-rugged sensitivity landscape.
Therefore, we set the non-convexity weighting factor $\lambda_{\text{rug}} = 0.0$, simplifying the test-time optimization objective to:
\begin{equation}
\label{eq:tta_loss}
\mathcal{L}_{\text{TTA}}(\theta_S) = \sum_{k \in S} \left[ \sum_{l=1}^L \mathcal{S}_k^{(l)}\left(\alpha_k(l)\right) \right]
\end{equation}
This formulation guarantees that any observed degradation in online methods is driven solely by transductive stream overfitting to noise, rather than artificial landscape ruggedness.
To give online TTA ample optimization capability, we extend its step budget to $T_{\text{steps}} = 100$ gradient updates using first-order Adam.

\subsubsection{Multi-Task Routing and the Privileged TTA Confound}
In real-world deployment scenarios, test-time inference often operates on a mixed, unlabeled streaming target distribution containing samples from multiple distinct tasks (e.g., a mixed stream of digits and clothing items).
During unsupervised online TTA, because the model has no prior knowledge of which task is active for a given input batch, it must minimize prediction entropy jointly across all task-specific heads.
This standard setting is called \emph{Unsupervised TTA}.
We explicitly qualify that this "privilege trap" and the resulting joint optimization challenge primarily manifest under heterogeneous, interleaved task streams.
In contrast, standard evaluation protocols in prior adaptive model-merging publications often evaluate on sequential single-task streams (e.g., a pure stream of CIFAR-10, followed sequentially by a pure stream of SVHN).
Under such sequential single-task conditions, the online optimizer adapts to only one task's active distribution at any given time, which naturally bypasses the need for joint multi-head entropy minimization.
However, we argue that assuming sequential single-task streams is an over-simplification of real-world deployments.
In practical multi-task environments, incoming streams are inherently interleaved and unlabeled, forcing the system to either perform joint multi-head entropy minimization (which triggers representation collapse) or rely on privileged, oracle task-routing labels (which violates the unsupervised assumption).
However, to prevent optimization collapse in high-conflict regimes, practitioners sometimes employ a simplified setting where ground-truth task labels are supplied at test-time to route the gradients through the active task head alone.
This is called \emph{Privileged TTA}.
While Privileged TTA stabilizes optimization, it represents a major, un-realistic deployment confound (the "privilege trap") by assuming the existence of oracle routing labels at inference time.
In Section~\ref{subsec:physical_validation}, we evaluate online TTA under both Unsupervised and Privileged settings, revealing that unconstrained methods collapse under both regimes when representation conflict is severe, whereas offline OFS-Tune operates without any privileged task routing assumptions during deployment.

To implement inference-time prediction routing in practice on interleaved mixed streams (where selecting the correct output head is required), developers can deploy simple and efficient strategies.
For example, one can utilize a zero-shot classifier (such as a vision-language model like CLIP) to identify the task domain, or train a lightweight, high-speed routing head (e.g., a simple logistic regression or single-layer MLP) on the same $M=10$ labeled validation samples per task.
Because this routing classifier is trained offline and contains negligible parameters, it can route test samples to their active heads at inference time with minimal latency, further demonstrating the practical deployment feasibility of OFS-Tune.

\subsubsection{Correlated Stream Noise Modeling}
A critical methodological flaw in prior adaptive merging simulators is the assumption of independent, zero-mean noise at each gradient update step.
In physical systems, online TTA adapts on a single local batch (or a highly correlated local stream), meaning stream-level selection bias is constant or highly correlated across adaptation steps.
To model this realistically, we sample a transductive stream noise offset $\epsilon_{\text{stream}} \sim \mathcal{N}(0, 0.10)$ \textbf{exactly once per adaptation session}, adding it directly to the target optimal profiles:
\begin{equation}
\label{eq:stream_noise}
\alpha_{k, \text{opt}}^{(l),\text{stream}} = \alpha_{k, \text{opt}}^{(l)} + \epsilon_{\text{stream}}
\end{equation}
This forces the online optimizer to navigate real-world transductive batch bias.
If the optimizer is over-parameterized (such as unconstrained AdaMerging), it will overfit to this stream-level noise, causing performance collapse.

\subsection{Offline Few-Shot Validation Tuning (OFS-Tune)}
In contrast to online TTA, Offline Few-Shot Validation Tuning (OFS-Tune) assumes that a tiny, labeled validation set containing $M=10$ samples per task is available offline.
The optimization objective is strictly supervised:
\begin{equation}
\label{eq:ofs_loss}
\min_{\theta_S} \sum_{k \in S} \left[ \sum_{l=1}^L \mathcal{S}_k^{(l)}\left(\alpha_k(l; \theta_S)\right) \right] + \epsilon_{\text{val}}
\end{equation}
where $\epsilon_{\text{val}} \sim \mathcal{N}(0, 0.01)$ represents the minor, independent variance introduced by sampling a few-shot validation set.
Crucially, random uniform draws of size $M=10$ from multi-class datasets (such as $C=10$ classes in CIFAR-10) introduce a severe risk of omitting classes entirely (a $99.96\%$ probability under naive sampling), causing extreme validation class imbalance.
To resolve this, we employ \emph{stratified sampling} to ensure that exactly $M / C$ samples are drawn from each class, guaranteeing uniform label coverage.
We provide a complete mathematical formulation of this class-imbalance risk and analyze label-space budget scaling limits in Appendix~\ref{subsec:class_imbalance}.
Crucially, because OFS-Tune combines a low-degree continuous polynomial constraint (such as linear $d=1$ or quadratic $d=2$) with a supervised objective, the optimization is highly regularized.
We optimize the parameters $\theta_S$ offline using Nelder-Mead derivative-free local search to minimize the validation loss.
While global derivative-free alternatives like Bayesian Optimization can be used, we find that Nelder-Mead is exceptionally robust and fast-converging for this low-dimensional search space (we provide a detailed empirical comparison of Nelder-Mead, Bayesian Optimization, and Random Search in Appendix~\ref{app:solver_config}).
Because the low-degree polynomial constraint rejects high-frequency noise, OFS-Tune acts as an analytical low-pass filter, completely ignoring validation sampling noise and generalizing robustly to the true underlying optimal profiles.

\paragraph{Setting and Data Access Trade-offs:}
We explicitly highlight that online TTA and offline OFS-Tune operate under fundamentally distinct assumptions regarding data access and deployment settings.
Online TTA is designed for fully unsupervised, zero-shot environments at test time, where labeled validation data is completely unavailable, or privacy and operational constraints strictly prohibit storing task-specific validation samples.
In such extreme settings, online unsupervised adaptation remains the only viable choice for dynamic coefficient adjustment.
In contrast, OFS-Tune represents a supervised, few-shot offline paradigm.
It leverages a tiny, labeled calibration set ($M=10$ samples per task) during a pre-deployment tuning phase.
While this offline assumption is highly realistic and common in practical engineering workflows (e.g., model staging and regression testing), it restricts OFS-Tune's applicability in strictly zero-shot, data-absent scenarios.
Thus, our work establishes a key trade-off for practitioners: online TTA offers zero-shot flexibility at the cost of vulnerability to transductive noise and potential representation collapse, whereas OFS-Tune offers robust noise filtering and zero test-time overhead at the cost of requiring a minimal labeled offline calibration set.

\subsection{Offline Unconstrained Validation Tuning (OFS-Unconstrained)}
To isolate the regularizing effect of the polynomial trajectory constraint from the effect of having supervised few-shot validation data, we introduce a crucial ablation baseline: **Offline Unconstrained Few-Shot Validation Tuning (OFS-Unconstrained)**.
OFS-Unconstrained utilizes the exact same labeled few-shot validation set ($M=10$ samples per task) and validation loss of Equation~\ref{eq:ofs_loss}, but optimizes $K \times L$ completely unconstrained layer-wise coefficients:
\begin{equation}
\label{eq:ofs_unc_loss}
\min_{\alpha_S} \sum_{k \in S} \left[ \sum_{l=1}^L \mathcal{S}_k^{(l)}\left(\alpha_k(l)\right) \right] + \epsilon_{\text{val}}
\end{equation}
where $\alpha_S \in \mathbb{R}^{K \times L}$.
We optimize these $K \times L$ parameters offline using L-BFGS-B with box constraints $[0, 1]$ to minimize the validation loss.
If few-shot validation data alone is sufficient for robust merging, then OFS-Unconstrained should perform equally well or better than OFS-Tune due to its higher capacity.
However, if unconstrained optimization overfits to validation set sampling variance $\epsilon_{\text{val}}$, then OFS-Unconstrained will exhibit degraded generalizability compared to the constrained OFS-Tune.
This baseline allows us to empirically isolate the regularizing value of the polynomial depth constraint.

\subsection{Summary of Method Characteristics}
To make the distinct data-access, compute, and deployment assumptions of each model-merging methodology immediately clear, we contrast Uniform merging, online TTA methods, and our proposed offline OFS-Tune framework across five key technical dimensions in Table~\ref{tab:method_comparison}.
We explicitly note that in environments where absolutely zero labeled validation data can be accessed (e.g., due to extreme user-privacy regulations or zero-shot constraints), OFS-Tune cannot be deployed as a drop-in replacement, and unsupervised online TTA remains the only viable choice despite its optimization risks.

\begin{table*}[t]
\caption{Systematic comparison of architectural and deployment assumptions across different model-merging methodologies.
OFS-Tune achieves zero test-time overhead and complete immunity to online stream noise by shifting parameter optimization to a brief, highly-regularized offline pre-deployment phase.}
\label{tab:method_comparison}
\vskip 0.15in
\begin{center}
\begin{scriptsize}
\setlength{\tabcolsep}{2.0pt}
\begin{tabular*}{\textwidth}{@{\extracolsep{\fill}}lccccc}
\toprule
\textbf{Method} & \textbf{Data Req.} & \textbf{Test Comp.} & \textbf{Oracle ID} & \textbf{Noise Risk} & \textbf{Phase} \\
\midrule
\textbf{Uniform} & None & None & No & None & None \\
\textbf{AdaMerging} & Unlabeled & Fwd-Bwd & Req.\textsuperscript{a} & Severe (Collapse) & Online (TTA) \\
\textbf{PolyMerge} & Unlabeled & Fwd-Bwd & Req.\textsuperscript{a} & High (Overfit) & Online (TTA) \\
\textbf{OFS-Tune (Ours)} & Few-shot ($M=10$) & None & No & None & Offline (Pre-deploy) \\
\bottomrule
\end{tabular*}
\end{scriptsize}
\end{center}
\vskip -0.1in
\begin{center}
\begin{minipage}{0.9\textwidth}
\textsuperscript{a}\scriptsize{~Online TTA requires oracle task labels at test-time (the "privilege trap") to route gradients through correct task heads; otherwise, joint unsupervised entropy minimization causes severe representation collapse.}
\end{minipage}
\end{center}
\end{table*}

\subsection{Methodological Takeaway: The Surrogate Loss Mismatch}
\label{subsec:methodological_takeaway}
We explicitly establish a critical methodological takeaway regarding the evaluation of model-merging algorithms.
In our simulation study, the online TTA optimization objective $\mathcal{L}_{\text{TTA}}$ (Equation 6) directly tracks the ground-truth optimal parameter profiles perturbed by noise.
This smooth, parameter-tracking surrogate loss represents a highly idealized and simplified setting.
In actual physical deployments, online TTA has zero access to optimal trajectories or task-specific sensitivity curves and must optimize a highly non-convex, rugged, and potentially misaligned unsupervised prediction entropy surface.
Consequently, abstracting high-dimensional weight merging via smooth, parameter-tracking surrogate losses in mathematical simulators inherently underestimates the instability, gradient noise, and catastrophic degradation risks of actual unsupervised entropy minimization.
We highlight this surrogate loss mismatch as a vital lesson for the model-merging community: simulation results should be treated as an optimistic upper bound, and rigorous physical weight-space validation is mandatory to expose the severe representation collapse risks inherent to actual unsupervised test-time optimization.
